\documentclass[11pt,a4paper]{article}

\usepackage[a4paper,margin=2.5cm]{geometry}
\usepackage[utf8]{inputenc}
\usepackage[T1]{fontenc}
\usepackage{amsmath,amssymb}
\usepackage{booktabs}
\usepackage{array}
\usepackage{graphicx}
\usepackage{xcolor}
\usepackage{hyperref}
\usepackage{listings}
\usepackage{enumitem}
\usepackage{titling}
\usepackage{fancyhdr}
\usepackage{parskip}

\hypersetup{
  colorlinks=true,
  linkcolor=blue!50!black,
  citecolor=blue!50!black,
  urlcolor=blue!50!black,
  pdftitle={PDDF Regularization Investigation -- sasfit\_inversion},
  pdfauthor={Joachim Kohlbrecher}
}

\lstset{
  basicstyle=\ttfamily\small,
  breaklines=true,
  frame=single,
  columns=flexible,
  backgroundcolor=\color{gray!8},
  rulecolor=\color{gray!40}
}

\pagestyle{fancy}
\fancyhf{}
\fancyhead[L]{PDDF Regularization Investigation}
\fancyhead[R]{\thepage}
\renewcommand{\headrulewidth}{0.3pt}

\newcolumntype{Y}{>{\centering\arraybackslash}X}

\title{\textbf{Regularized Inversion of the PDDF Fredholm Integral in\\
\texttt{sasfit\_inversion}: Diagnosis, Elimination, and a\\
Hann-Tapered-Boundary Fix}}
\author{Joachim Kohlbrecher \\ \small Paul Scherrer Institut (PSI)}
\date{\today}

\begin{document}

\maketitle

\begin{abstract}
\noindent
The \texttt{sasfit\_inversion} package inverts small-angle-scattering
intensity $I(q)$ for the pair-distance-distribution function (PDDF)
$p(r)$ by solving a Fredholm integral equation of the first kind with
the signed $4\pi\, j_0(qr)$ kernel. Early regularized solvers on real
test data (\texttt{tests/data/test.dat}) produced a deceptively
acceptable aggregate $\chi^2_r$ that concealed a severe, high-$q$
noise-tracking pathology. This report documents the full
investigation: adoption of Hansen's (2000) boundary-constrained
Bayesian-evidence inverse Fourier transform (IFT) to fix a
rank-deficient smoothness operator; a region-resolved diagnostic
methodology that exposed the hidden high-$q$ overfitting; four
independent systematic scans that ruled out regularization strength,
$D_{max}$, grid density, and the inner boundary radius as the cause;
a two-hyperparameter adaptive-$\lambda$ experiment that failed for a
principled reason (the oscillatory kernel's broadband sensitivity to
all of $r$); an explicitly rejected non-negativity constraint (invalid
for core--shell and contrast-variation particles); and the eventual
fix -- replacing Hansen's hard $p(D_{max})=0$ knot with a
Hann-tapered soft boundary -- validated both on synthetic data with
known ground truth and on the real pathological dataset, where it
improved the high-$q$ regional $\chi^2$ contribution by a factor of
nine (0.023 $\to$ 0.208) while leaving the already-reasonable
low-$q$ fit unchanged.
\end{abstract}

\tableofcontents

%=====================================================================
\section{The inversion problem}
%=====================================================================

Recovering the pair-distance-distribution function $p(r)$ from
small-angle scattering intensity $I(q)$ requires solving a Fredholm
integral equation of the first kind,
\begin{equation}
  I(q) \;=\; 4\pi \int_0^{D_{max}} p(r)\, \frac{\sin(qr)}{qr}\, dr ,
  \label{eq:fredholm}
\end{equation}
subject to the physical boundary conditions $p(0) = 0$ and $p(r) = 0$
for $r > D_{max}$, where $D_{max}$ is the particle's maximum linear
dimension. Equation~\eqref{eq:fredholm} is discretized on a point grid
in $r$ (direct grid estimation, superseding Glatter's original
small-basis-function restriction) and the resulting linear system
$\mathbf{K}\mathbf{p} = \mathbf{I}$ is classically ill-posed: the
kernel matrix $\mathbf{K}$ is severely ill-conditioned, so small
amounts of noise in the measured $I(q)$ map onto large, unphysical
oscillations in the naively recovered $p(r)$ unless the solution is
regularized.

An initial attempt used the signed $4\pi\, j_0(qr)$ kernel together
with a naive second-derivative smoothness penalty. This produced
catastrophically poor fits (reduced chi-square $\chi^2_r$ in the
hundreds), which traced to two compounding problems:
\begin{itemize}[nosep]
  \item the plain second-derivative operator has a two-dimensional
    null space -- a constant offset and a linear ramp in $p(r)$ are
    both invisible to the roughness penalty and therefore completely
    unconstrained;
  \item the kernel's own condition number for this dataset's $q$/$r$
    grid is of order $10^{16}$.
\end{itemize}

%=====================================================================
\section{Hansen's boundary-constrained Bayesian-evidence IFT}
%=====================================================================

Hansen (2000, \emph{J.\ Appl.\ Cryst.} \textbf{33}, 1415--1421) resolves
the null-space problem by building the physical boundary condition
$p(0) = p(D_{max}) = 0$ \emph{directly into the smoothness functional}
(his eq.~19), rather than imposing it through a restricted basis set:
\begin{equation}
  S(p) \;=\; \frac{1}{2}\sum_{i=0}^{N-2} (p_i - p_{i+1})^2
             \;+\; \frac{1}{2}p_0^2 \;+\; \frac{1}{2}p_{N-1}^2 .
  \label{eq:hansen-penalty}
\end{equation}
This operator has a provably full-rank Hessian (asserted at
construction time in
\texttt{regularization.hansen\_smoothness\_cholesky}), which drops the
problem's condition number from $\sim\!10^{16}$ to $\sim\!9.2\times
10^{3}$ on the real dataset.

The regularization strength $\lambda$ is chosen by \textbf{Bayesian
evidence maximization} (Vestergaard \& Hansen, 2006, \emph{J.\ Appl.\
Cryst.} \textbf{39}, 797--804) rather than generalized cross-validation
(GCV), the L-curve, or the discrepancy principle. The evidence
machinery is implemented generically in
\texttt{solvers.bayesian\_evidence} so that any regularization operator
$\mathbf{L}$ -- not only Hansen's -- can be plugged into the same
$\lambda$-selection procedure.

This combination is the \texttt{bayesian\_evidence\_hansen} solver.
On real data it was a dramatic improvement over every regularized
solver tried previously: $\chi^2_r \approx 1.6$, versus $\approx 586$
for EM with smoothing and $\approx 34{,}600$ for plain
Tikhonov regularization with GCV-selected $\lambda$, at the time this
solver was introduced.

%=====================================================================
\section{A hidden problem behind a good aggregate $\chi^2_r$: regional overfitting}
\label{sec:regional}
%=====================================================================

A good \emph{aggregate} $\chi^2_r$ can conceal a badly imbalanced fit.
Splitting the fit residuals $(I_{fit}(q) - I_{data}(q))/\delta I(q)$
into low-, mid-, and high-$q$ terciles and computing the mean squared
residual (chi-square contribution) per region revealed that the
Hansen solver's overall $\chi^2_r \approx 1.6$ was in fact a
compromise between two failures in opposite directions:
\begin{itemize}[nosep]
  \item \textbf{low-$q$ was underfit}, with a chi-square contribution
    of order 3 or more;
  \item \textbf{high-$q$ was severely overfit}, tracking the noise
    almost exactly (chi-square contribution $\approx 0.02$ -- i.e.\
    the fit matched the data roughly seven times more closely than the
    stated error bars justify).
\end{itemize}

A single global $\lambda$ cannot, in general, simultaneously be
correct for a region containing real structural information (which
wants less smoothing) and a region with little information content
(which wants more smoothing, to avoid chasing noise). The remainder of
the investigation asks \emph{why} this happens here, and whether any
choice of hyperparameter or grid can fix it. This region-resolved
diagnostic -- computing the mean $\chi^2$ contribution separately for
low-, mid-, and high-$q$ terciles -- is reused as the evaluation
methodology for every scan and experiment described below.

%=====================================================================
\section{Systematically ruling out the usual suspects}
\label{sec:scans}
%=====================================================================

Four independent scans were run on the real dataset, each holding
every other parameter fixed (scripts
\texttt{tests/diagnose\_hansen\_*.py}).

\subsection{$\lambda$ scan}

Scanning $\lambda$ across six decades showed that the high-$q$
chi-square contribution approaches 1 only once $\lambda$ is so large
that low-$q$ is \emph{already} badly underfit ($\chi^2 \gg 1$): there
is no single value of $\lambda$ that brings low-, mid-, and high-$q$
all close to 1 simultaneously. Critically, \textbf{low-$q$ chi-square
never dropped below $\approx 2.16$, even at the smallest $\lambda$
tested} (essentially unregularized, with 49 of 50 parameters
effectively free). More freedom does not fix low-$q$ -- the first
indication that this was not simply a regularization-strength
problem.

\subsection{$D_{max}$ scan}

$D_{max} = 500$\,\AA\ turned out to be genuinely too small, truncating
real pair-distance density: low-$q$ chi-square dropped from 3.35 to
$\approx 0.8$ as $D_{max}$ was increased from 500 to $\approx 800$.
This motivated checking Glatter's consistency condition,
\begin{equation}
  q_{min} \;\le\; \frac{\pi}{D_{max}}
  \quad\Longleftrightarrow\quad
  D_{max} \;\le\; \frac{\pi}{q_{min}} ,
  \label{eq:glatter}
\end{equation}
which for this dataset's $q_{min} = 0.004016\,\text{\AA}^{-1}$ gives
the largest physically constrainable value, $D_{max} \approx 782$\,\AA.
The GUI now auto-suggests this value on data load (overridable by the
user).

Beyond the Glatter bound, apparent improvement in the aggregate
$\chi^2_r$ was shown to be mostly overfitting (mid-$q$ chi-square kept
sliding well below~1) rather than genuinely improved resolution --
and, crucially, \textbf{the high-$q$ chi-square contribution stayed
pinned at $\approx 0.02$ across every $D_{max}$ tested}, confirming
that $D_{max}$ is not the mechanism either.

\subsection{Grid-density ($n_r$) scan}

This scan tested whether oversampling $r$ -- more grid points than the
data's Shannon-limited information content supports -- was handing the
high-$q$ region cheap, extra degrees of freedom with which to fit
noise. This was refuted: once above the minimum resolved grid density
($n_r \gtrsim 75$ for this problem), every metric, including the
high-$q$ chi-square contribution, was flat from $n_r = 75$ to
$n_r = 200$.

\subsection{Inner-boundary-radius ($r_{min}$) scan}

This scan tested the boundary-condition approximation (the constraint
$p(0)=0$ is in practice enforced at the grid's first node, $r_{min}$,
not exactly at $r = 0$) together with possible near-$r=0$ kernel
degeneracy. Also flat, from $r_{min} = 0.1$ to $r_{min} = 10$\,\AA\ --
no effect.

\medskip
\noindent\textbf{Conclusion.} $\lambda$, $D_{max}$, $n_r$, and
$r_{min}$ were each ruled out individually. The high-$q$
noise-tracking survived every one of them, unchanged, precisely in the
regime where the method is otherwise well-resolved.

%=====================================================================
\section{A failed fix: position-dependent (r-segment) adaptive $\lambda$}
\label{sec:adaptive}
%=====================================================================

Reasoning that a single global $\lambda$ was the limiting factor,
Hansen's penalty~\eqref{eq:hansen-penalty} was split into two additive
pieces by $r$-index,
\begin{equation}
  \mathbf{A}_{\text{Hansen}} \;=\; \mathbf{A}_{lo} + \mathbf{A}_{hi} ,
\end{equation}
an exact algebraic partition implemented in
\texttt{regularization.hansen\_smoothness\_two\_segment}, with an
independent $\lambda_{lo}$, $\lambda_{hi}$ chosen by a generalized
two-hyperparameter evidence search
(\texttt{bayesian\_evidence.log\_evidence\_blocks} /
\texttt{evidence\_search\_blocks}), the block formulation chosen so
that it collapses exactly to the single-$\lambda$ evidence formula
when $\lambda_{lo} = \lambda_{hi}$. Both the split location and the
two $\lambda$ values were searched jointly.

\textbf{Result: no improvement.} The high-$q$ chi-square contribution
was unchanged or slightly worse (0.0234 $\to$ 0.0298) regardless of
where the split was placed, and $\lambda_{lo}$ pinned at the search
grid's floor for every candidate split location.

The reason, in hindsight, is structural rather than a search-grid
artifact: the oscillatory $j_0(qr)$ kernel means a single high-$q$
data point is sensitive to $p(r)$ across the \emph{whole} $r$-range,
not a contiguous region -- so a noise-fitting wiggle in the recovered
$p(r)$ can sit anywhere in $r$. Partitioning the regularization
penalty by $r$-\emph{position} does not target the actual mechanism,
which turned out to be spatial-frequency content, not position.

%=====================================================================
\section{A rejected fix: non-negativity}
\label{sec:nonneg}
%=====================================================================

$p(r) \ge 0$ is \emph{not} a universal constraint for this package's
intended scope -- it holds only for a single, homogeneous-contrast
particle. For core--shell particles, multi-shell/multilamellar
structures, or contrast-variation (partially deuterated) systems, the
excess scattering-length density changes sign within the particle, and
$p(r)$ can legitimately dip negative from cross-correlation between
oppositely-signed regions. This is a documented diagnostic signature
of internal structure in the SAXS/SANS literature, not a numerical
artifact.

Imposing $p(r) \ge 0$ as a hard, default constraint would distort or
suppress genuine structural signal for exactly this class of system,
which lies within the package's intended scope. \textbf{It is not
implemented as a default}; if offered at all, it must be an explicit
opt-in variant restricted to the homogeneous-contrast case.

%=====================================================================
\section{The fix: a Hann-tapered boundary}
\label{sec:hann}
%=====================================================================

$p(r)$ and $I(q)$ are, in the relevant sense, a Fourier sine-transform
pair. Hansen's boundary condition is mathematically a \emph{sharp
edge} in $r$-space: $p$ is pinned to exactly zero at a single grid
point. By the standard windowing/Gibbs-phenomenon relationship, a
sharp edge in one domain convolves the conjugate domain's transform
with a sinc function, whose sidelobes decay slowly ($\sim 1/q$) and
persist out to high $q$ -- \emph{wherever} the edge is placed. This
single mechanism explains both negative results above: the $D_{max}$
scan (Section~\ref{sec:scans}) does not fix the pathology because
moving the edge does not remove it, and $r$-segment splitting
(Section~\ref{sec:adaptive}) does not fix it because the resulting
ringing is not localized to one region of $r$.

The remedy is the standard apodization fix for truncation ringing:
replace the hard edge with a smooth rolloff. The function
\texttt{regularization.hann\_taper(r, r\_max, taper\_frac=0.25)}
returns a weight $w(r) \in [0,1]$, equal to 1 in the interior and
tapering to 0 with a raised-cosine (Hann-style) profile over the last
quarter of $[0, D_{max}]$:
\begin{equation}
  w(r) =
  \begin{cases}
    1, & r \le (1-\tau)\,D_{max} \\[4pt]
    \dfrac{1}{2}\Bigl(1 + \cos\!\bigl(\pi\, x(r)\bigr)\Bigr), &
      (1-\tau)\,D_{max} < r \le D_{max} \\[8pt]
    0, & r > D_{max}
  \end{cases}
  \qquad
  x(r) = \frac{r - (1-\tau)D_{max}}{\tau\, D_{max}} ,
  \label{eq:hann-taper}
\end{equation}
with taper fraction $\tau = \texttt{taper\_frac} = 0.25$. Reparametrizing
\begin{equation}
  p(r) \;=\; w(r)\, p_{free}(r)
\end{equation}
and solving for $p_{free}$ with a \emph{plain} (non-boundary-constrained)
second-derivative smoothness penalty replaces Hansen's hard knot with
a soft decay. The plain second-derivative operator's null space --
harmless when combined with Hansen's own boundary fix -- needed a
tiny numerical ridge stabilizer in this variant, since nothing else
constrains it (the \texttt{ridge} parameter threaded through
\texttt{bayesian\_evidence.py}'s \texttt{\_map\_solution},
\texttt{log\_evidence}, and \texttt{evidence\_search}).

\subsection{Validation on synthetic data (known ground truth)}

A homogeneous-sphere PDDF ($D_{true} = 380$\,\AA, with the standard
analytic form $p(r) = r^2\bigl(1 - \tfrac{3}{2}\tfrac{r}{D} +
\tfrac{1}{2}(\tfrac{r}{D})^3\bigr)$ for $0 \le r \le D$) was
forward-transformed to $I(q)$ on a $q$-grid matching the real
dataset's span and point count, with realistic relative-plus-floor
noise added, across five independent noise draws. Table~\ref{tab:synth}
compares the recovered $p(r)$'s RMSE against the known true $p(r)$,
and the high-$q$ chi-square contribution, for the hard (Hansen)
boundary versus the Hann-tapered boundary.

\begin{table}[h]
\centering
\caption{Synthetic sphere validation, 5 independent noise seeds.
Lower RMSE and a high-$q$ chi-square contribution closer to 1 are both
better.}
\label{tab:synth}
\begin{tabular}{@{}c c c c c@{}}
\toprule
 & \multicolumn{2}{c}{RMSE vs.\ true $p(r)$} & \multicolumn{2}{c}{high-$q$ $\chi^2$ contribution} \\
\cmidrule(lr){2-3}\cmidrule(lr){4-5}
seed & hard boundary & Hann-tapered & hard & tapered \\
\midrule
1 & 0.0474 & 0.0409 & 0.614 & 0.633 \\
2 & 0.0381 & 0.0233 & 0.839 & 0.885 \\
3 & 0.0434 & 0.0266 & 1.141 & 1.196 \\
4 & 0.0362 & 0.0295 & 1.192 & 1.260 \\
5 & 0.0566 & 0.0240 & 0.477 & 0.580 \\
\bottomrule
\end{tabular}
\end{table}

The Hann-tapered version had lower RMSE against the known true $p(r)$
in \emph{every single trial} (by roughly 35--58\% in several cases),
and a high-$q$ chi-square contribution consistently, if modestly,
closer to 1 (i.e.\ less overfitting) in every trial as well. This
synthetic test -- a smooth, textbook sphere PDDF -- is inherently
better-conditioned than the real particle's actual shape, so it
validates the \emph{mechanism and its direction}, not the full
\emph{severity} seen on real data.

\subsection{Validation on real data}

Table~\ref{tab:real} shows the region-resolved chi-square breakdown on
the actual pathological dataset
(\texttt{tests/diagnose\_hansen\_tapered.py}), comparing the
\texttt{bayesian\_evidence\_hansen} (hard boundary) and
\texttt{bayesian\_evidence\_hann\_tapered} solvers.

\begin{table}[h]
\centering
\caption{Region-resolved $\chi^2$ contribution on the real test
dataset. The ideal value for a correctly weighted region is 1.}
\label{tab:real}
\begin{tabular}{@{}lcc@{}}
\toprule
region & hard boundary (\texttt{bayesian\_evidence\_hansen}) &
  Hann-tapered (\texttt{bayesian\_evidence\_hann\_tapered}) \\
\midrule
low-$q$ $\chi^2$  & 0.797 & 0.807 \\
mid-$q$ $\chi^2$  & 0.343 & 0.491 \\
\textbf{high-$q$ $\chi^2$} & \textbf{0.023} & \textbf{0.208} \\
\midrule
overall $\chi^2_r$ & 0.384 & 0.499 \\
\bottomrule
\end{tabular}
\end{table}

On the actual pathological dataset, the high-$q$ chi-square
contribution moved \textbf{nine times} closer to the ideal value of 1
($0.023 \to 0.208$), mid-$q$ also improved ($0.343 \to 0.491$), and
low-$q$ -- which was already reasonable -- was essentially unchanged.
This is a genuine, substantial improvement in regional balance, not
merely the synthetic test's direction holding up in principle.

It is \textbf{not a complete fix}: a high-$q$ chi-square contribution
of 0.208 is still well below the ideal value of 1, so some residual
overfitting remains. The taper fraction $\tau$ (currently a fixed,
untuned 0.25) has not been scanned for an optimum, and doing so is a
natural next step for anyone picking this work back up
(Section~\ref{sec:future}).

%=====================================================================
\section{Current recommendation}
%=====================================================================

For the $j_0(qr)$/PDDF kernel on real data:
\begin{itemize}[nosep]
  \item Use \texttt{bayesian\_evidence\_hann\_tapered} as the primary
    choice -- it consistently gives better regional balance than the
    hard-boundary version, validated on both synthetic (known ground
    truth) and real data.
  \item \texttt{bayesian\_evidence\_hansen} (hard boundary) remains
    available, in particular as a simpler/more classical baseline to
    compare against, or for cases where the taper's extra
    approximation (the ridge stabilizer) is undesirable.
  \item Both are registered in
    \texttt{sasfit\_inversion.solver\_registry.SOLVER\_REGISTRY} and
    therefore automatically available in the GUI's solver dropdown,
    with their full trade-offs documented in each
    \texttt{SolverSpec.description}.
  \item For the sphere/size-distribution kernel (non-negative, not
    PDDF), EM with smoothing and discrepancy-principle $\lambda$
    selection remains the better-validated default -- none of this
    investigation concerned that kernel.
\end{itemize}

%=====================================================================
\section{Open questions for future work}
\label{sec:future}
%=====================================================================

\begin{itemize}[nosep]
  \item \textbf{Tune \texttt{taper\_frac}.} Not yet scanned for an
    optimum; the current value of 0.25 is a reasonable but untuned
    default.
  \item \textbf{Check high-$q$ error-bar calibration.} Compare raw
    point-to-point $I(q)$ scatter against the quoted $\delta I$,
    independent of any inversion. If $\delta I$ is overestimated at
    high $q$, part of the apparent ``overfitting'' signature may be
    the solver correctly reporting that the data contain less noise
    than their stated error bars claim, rather than a regularization
    failure. Not yet checked.
  \item \textbf{Frequency-domain (rather than position-domain) adaptive
    regularization.} E.g.\ a higher-order roughness penalty targeting
    short-wavelength ripples specifically, wherever they occur in $r$,
    was proposed but not implemented; it may compound usefully with
    the Hann taper rather than replace it.
  \item \textbf{Total-variation (L1-on-the-gradient) regularization}
    as an alternative to Hansen's L2 smoothness was proposed
    (sparsity-promoting, suppresses ringing without a sign assumption,
    unlike non-negativity) but not implemented -- it would require
    moving from a closed-form quadratic solve to an IRLS/QP
    formulation.
\end{itemize}

%=====================================================================
\section*{References}
\addcontentsline{toc}{section}{References}
%=====================================================================
\begin{itemize}[nosep]
  \item Hansen, S. (2000). \emph{Bayesian estimation of hyperparameters
    for indirect Fourier transformation of small-angle scattering
    data.} J.\ Appl.\ Cryst.\ \textbf{33}, 1415--1421.
  \item Vestergaard, B. \& Hansen, S. (2006). \emph{Application of
    Bayesian analysis to indirect Fourier transformation in
    small-angle scattering.} J.\ Appl.\ Cryst.\ \textbf{39}, 797--804.
  \item Glatter, O. (1977). \emph{A new method for the evaluation of
    small-angle scattering data.} J.\ Appl.\ Cryst.\ \textbf{10},
    415--421.
\end{itemize}

\end{document}
